\chapter*{Abstract}
\addcontentsline{toc}{chapter}{Abstract}
This thesis presents the design, implementation, and evaluation of a low-power LoRa-based wireless sensor network for monitoring soil and environmental conditions beneath Icelandic turf roofs. The system consists of sensor nodes built around Arduino Pro Mini microcontrollers, with capacitive soil moisture sensors, DS18B20 soil temperature sensors, AHT10 air temperature and humidity sensors, local SD card storage, and 433 MHz LoRa communication. A base station forwards collected measurements to the ThingSpeak cloud platform for remote access and visualization.

Laboratory and field tests were conducted to evaluate power consumption, communication performance, sensor calibration, and cloud data transmission. Under the defined operating cycle, the sensor node consumed an average current of approximately 3.08 mA. With three 3,500 mAh battery cells connected in parallel, providing a nominal capacity of 10,500 mAh, the estimated operating lifetime was approximately 142 days. This is an idealized estimate and does not account for reductions in usable battery capacity at low temperatures.

Field testing in the semi-urban area of Vellir in Hafnarfj\"{o}r\dh ur observed successful single-hop LoRa communication at several test locations within an approximate distance of 3 km using 433 MHz antennas, with recorded RSSI values from -103 dBm to -113 dBm. Communication performance varied with antenna placement and environmental obstacles, and occasional packet loss was observed in obstructed areas. Calibration of the capacitive soil moisture sensor produced a second-order relationship over the tested 0--40\% moisture range, with an in-sample root mean square error of approximately 3.01 percentage points against the gravimetric calibration targets.

During cloud integration tests, ThingSpeak uploads took approximately 1--2 seconds on average. \draftchange{The system was also deployed at Keldur at Rangárvellir to monitor environmental conditions beneath four turf roofs,} with first-to-last valid record spans of approximately 51 to 130 days across the four sensor nodes. The results demonstrate the feasibility of implementing a low-cost monitoring system for Icelandic turf roofs, while highlighting the need to consider environmental conditions and long-term hardware reliability.

\vfill
\textbf{Keywords:} wireless sensor network (WSN), LoRa, mesh network, soil moisture, turf roofs, environmental monitoring, low-power systems, ThingSpeak

